%! Two Categorical Variables — Unit 2, Lesson 2.1 — Lesson Plan
% GRADUAL-RELEASE plan (warm-up / guided notes and practice / debrief / close and start homework).
% THERE IS NO GROUP ACTIVITY IN THIS COURSE — the guided notes carry the release: I do / we do
% row by row, then Guided Practice together. THERE IS NO INDEPENDENT PRACTICE SET — the homework
% is the individual practice, and the period ends with students starting it. Teacher-facing.
% Every box is `skillbox` (breakable) with a \boxguard ahead of it; this course has no
% `fixedskillbox`, no tier boxes, and no exit ticket.
% COMPACT BUDGET (2026-09-29): the packet is cover 1 · warm-up 1 · notes 2 · AP practice 1 ·
% homework 1, finished in one 55-minute period. Unit 2 folds CED Topics 2.1–2.3 into this lesson.
\documentclass[10pt]{article}
\usepackage{apstats-article}
\usepackage{apstats-boxes}

\newcommand{\UnitNumberName}{Unit 2: Exploring Two-Variable Data \quad}
\newcommand{\LessonNumberName}{Lesson 2.1: Two Categorical Variables}

\begin{document}
\begin{center}
    {\Large\bfseries \CourseName \\
    {\small\bfseries \UnitNumberName \LessonNumberName}}
\end{center}

\begin{tcolorbox}[colback=sky, colframe=navy, boxrule=0.9pt, arc=2mm,
    left=2mm, right=2mm, top=1.5mm, bottom=1.5mm]
    \textbf{Primary Objective:} Students will read a two-way table of two categorical variables,
    calculate joint, marginal, and conditional relative frequencies with the correct
    denominator, and use conditional distributions (in a table or a segmented bar graph) to
    decide whether the variables are associated --- stated in context, without claiming cause.
    \\[2pt]
    \textbf{CED alignment:} Topics 2.1, 2.2, 2.3 \quad$\cdot$\quad Big Ideas: VAR (Variation and
    Distribution), UNC (Patterns and Uncertainty) \quad$\cdot$\quad Skills 1.A, 2.C, 2.D
    \quad$\cdot$\quad LO VAR-1.D, UNC-1.P, UNC-1.Q, UNC-1.R \quad$\cdot$\quad EK VAR-1.D.1,
    UNC-1.P.1--1.P.4, UNC-1.Q.1, 1.Q.2, UNC-1.R.1
    \\[2pt]
    \textbf{Lesson model:} \emph{Gradual release --- I do, we do, you do --- all of it inside
    the guided notes.} There is no group activity. The notes name the vocabulary as they teach
    it and deliver the instruction row by row in one \emph{Main Ideas / Notes} table --- each
    idea a definition to complete and a grid of short problems, the first worked by you and the
    rest by the class; the last row, Guided Practice, is worked together with students holding
    the pen; \textbf{the homework is the individual practice}, and students start it in class
    while you circulate. The whole-class debrief is \emph{spoken}.
    \\[2pt]
    \textbf{Note on the ``no sketch'' rule:} the segmented bar graph is \textbf{given}. Students
    construct nothing but fractions from a table; they read the graph and argue from it.
\end{tcolorbox}

\renewcommand{\arraystretch}{1.4}
\boxguard
\begin{skillbox}[Learning Targets \& Key Understandings]{goldbox}
\begin{tabularx}{\linewidth}{@{}X|X@{}}
\textbf{Learning targets (student language)} & \textbf{Key understandings (the ``why'')} \\[2pt]
\hline
\vspace{-6pt}
\begin{itemize}[leftmargin=1.1em, nosep, after=\vspace{2pt}]
  \item I can ask whether two variables are related, and read a two-way table.
  \item I can find joint, marginal, and conditional relative frequencies, choosing the right
        total to divide by.
  \item I can read a segmented bar graph and say whether two variables are associated,
        with percents, in context.
  \item I can explain why an association in a survey does not show cause.
\end{itemize}
&
\vspace{-6pt}
\begin{itemize}[leftmargin=1.1em, nosep, after=\vspace{2pt}]
  \item A two-way table summarizes two categorical variables on the same individuals; a
        joint relative frequency is a cell over the grand total, a marginal one a row or
        column total over the grand total (EK UNC-1.P.3, 1.P.4, 1.Q.1).
  \item A conditional relative frequency restricts to one row or column and divides by that
        group's total (EK UNC-1.Q.2). Segmented, side-by-side, and mosaic bar graphs draw
        these conditional distributions (EK UNC-1.P.1).
  \item Variables are associated when the conditional distributions differ across groups
        (EK UNC-1.P.2, 1.R.1); an apparent pattern may be random or confounded (EK VAR-1.D.1).
  \item \textbf{Target misconception:} judging a relationship from \emph{counts} or the
        \emph{wrong total}. ``48 phone-in-room students sleep 8+ hours vs.\ only 20 without''
        compares groups of 160 and 40; the conditional rates are 30\% vs.\ 50\%, the
        \emph{opposite} story. Close cousin: dividing by the grand total (joint) or by the
        wrong group's total when a conditional was asked.
\end{itemize}
\end{tabularx}
\end{skillbox}

\boxguard
\begin{skillbox}[{Vocabulary, Concepts, \& Theorems --- taught directly in the guided notes}]{greenbox}
\begin{minipage}[t]{0.48\linewidth}
    \begin{itemize}
        \item \textbf{Two-way (contingency) table} --- counts for two categorical variables
              (EK UNC-1.P.3)
        \item \textbf{Joint relative frequency} --- cell $\div$ grand total (EK UNC-1.P.4)
        \item \textbf{Marginal relative frequency} --- row or column total $\div$ grand total
              (EK UNC-1.Q.1)
        \item \textbf{Conditional relative frequency} --- cell $\div$ its own row or column
              total (EK UNC-1.Q.2)
    \end{itemize}
\end{minipage}
\hfill
\begin{minipage}[t]{0.48\linewidth}
    \begin{itemize}
        \item \textbf{Segmented bar graph} --- one 100\% bar per group; \textbf{mosaic plot}
              --- bar widths show group sizes; \textbf{side-by-side bar graph} --- named
              aloud (EK UNC-1.P.1)
        \item \textbf{Association} --- conditional distributions differ across groups
              (EK UNC-1.P.2, 1.R.1)
        \item \textbf{Carried forward} --- categorical vs.\ quantitative (1.1); relative
              frequency and bar graphs (1.2)
    \end{itemize}
\end{minipage}
\end{skillbox}

% A tabularx never splits, so guard generously ahead of it.
\boxguard[30]
\begin{skillbox}[Lesson at a Glance --- \MeetingLength]{sky}
\renewcommand{\arraystretch}{1.25}
\begin{tabularx}{\linewidth}{@{}l c X X@{}}
\textbf{Phase} & \textbf{Min} & \textbf{Students} & \textbf{Teacher} \\
\hline
Warm-Up & 5 & A relative frequency, two quiz sections, and a ``related?'' question & Debrief item~2 aloud; leave \emph{more} vs.\ \emph{better} on the board all period \\
Guided Notes \& Practice & 34 & Fill the vocabulary box and the notes table: problems 1--10 row by row, then \emph{Texting Drivers} (11--14) together, pen in their hand & \textbf{I do} each row's definition lines and first problem (1, 4, 7), \textbf{we do} the rest of its grid (24) $\to$ \textbf{we do} Guided Practice (10) \\
Debrief & 8 & Pens down, papers up --- answer aloud, nothing to fill in & Open on the almost-right problem~8 answer; set counts beside percents for both contexts \\
Close \& Assign & 8 & \textbf{Start the homework in class}, alone and silent & Announce the paper homework, circulate for your formative read on A2 and B2, preview Lesson~2.2 \\
\end{tabularx}
\end{skillbox}

\boxguard[20]
\begin{skillbox}[Warm-Up --- Activate Prior Knowledge (5 min)]{sky}
\begin{minipage}[t]{0.48\linewidth}
    \textbf{The three items and what each seeds}
    \begin{itemize}
        \item \textbf{1.} 45 of 180 walk: a relative frequency. \textbf{Spirals} to Lesson 1.2
              and is the arithmetic of every joint and marginal in row~1.
        \item \textbf{2.} Section A 12 of 20, Section B 18 of 40. \textbf{Seeds the crux}:
              B had \emph{more} pass, A did \emph{better}. That is problem~8 in miniature.
        \item \textbf{3.} Classify three variables and ask whether two are related.
              \textbf{Seeds} Topic~2.1 (VAR-1.D) --- and it is the notes' own context.
    \end{itemize}
\end{minipage}
\hfill
\begin{minipage}[t]{0.48\linewidth}
    \textbf{Running it}
    \begin{itemize}
        \item \textbf{Debrief item~2 aloud} and write the two answers side by side:
              \emph{more pass: B} \quad \emph{did better: A}. Leave them up; in row~3 you
              point back at them.
        \item Take one student question from item~3 --- ideally ``do students with the phone
              in the room sleep less?'' --- and tell them that is exactly the question the
              notes answer.
        \item Item~1 is 45 seconds. Do not let it expand.
    \end{itemize}
\end{minipage}
\end{skillbox}

\boxguard
\begin{skillbox}[Guided Notes \& Practice --- I do, then we do (34 min)]{sky}
\textbf{This is the whole lesson.} There is no group activity, and there is no independent
practice set on this page: \textbf{the homework is the individual practice}, started in class.
The notes are one \emph{Main Ideas / Notes} table --- four rows, 14 short problems, on one data
set for all three instruction rows (200 students: phone in the bedroom $\times$ 8+ hours of
sleep). \textbf{In every row you fill the definition lines and work the first problem of the
grid; the class works the rest, pen in their hand.} The vocabulary box (three terms) is filled
\emph{as you go}, never in advance.
\begin{multicols}{2}
\textbf{Two-Way Table (6 min).} Name the table, then the two ways to divide by the grand
total. \textbf{You work problem~1} (joint, $48/200$); problems 2--3 are theirs. Problem~3 is
classification, not arithmetic: \emph{0.80 keep a phone in the room} is a row total over 200 ---
marginal. Fill \emph{two-way table} in the vocabulary box here.

\textbf{Conditional (8 min) --- the trap.} The rule to land: \emph{the words after ``of'' name
the denominator.} \textbf{You work problem~4} ($48/160$); they do 5. \textbf{Problem~6 is the
trap}: ``of 8+ hour sleepers'' flips the group to the \emph{column}, so it is $48/68 \approx 0.71$,
not $0.30$ and not $48/200$. Take a hand count on the denominator before anyone computes. Fill
\emph{conditional relative frequency}.

\textbf{Association (10 min) --- the crux.} The segmented bar graph is problem~4 and~5 drawn:
30\% vs.\ 50\%. Name the mosaic plot in one sentence (bar widths = group sizes) and say aloud
that a side-by-side bar graph draws the same percents as paired bars. \textbf{You work
problem~7} (read 70\%). \textbf{Problem~8, in this last instruction row, is the crux}: the claim
``48 phone students sleep 8+ vs.\ only 20'' is true in counts and backwards in rates. Point at
the warm-up's \emph{more} vs.\ \emph{better}. Problem~9 states the association with percents;
problem~10 is the hedge --- a survey, so no cause (and the pattern could even be chance,
VAR-1.D.1). Fill \emph{association}.

\columnbreak
\textbf{We do (about 10 min): Texting Drivers.} Problems 11--14, a fresh context, worked
entirely together. Students hold the pen; you hold the questions: \emph{``Of whom? Say the group
before you divide.''} \quad \emph{``Which group is bigger --- and does that make its count
bigger?''} \textbf{Problem~13 is the crux transferred} (36 older vs.\ 20 young texters, from
groups of 150 and 50) and is the one to protect when time is short; problem~14 must end in
context with the survey hedge.

\textbf{The pen must actually change hands.} Guided Practice is the only place today where you
can watch a student decide, so do not narrate it. Put problem~12 on them cold and take a hand
count on the two denominators (50 and 150) before anyone speaks.

Circulate and demand the reason, not the number:
\begin{itemize}[leftmargin=1.1em, nosep]
  \item \textbf{Problem~11:} ``Is that one cell or a whole row? Then which kind is it?''
  \item \textbf{Problem~12:} ``Read me your bottom number. Whose total is it?''
  \item \textbf{Problem~13, the crux:} ``Out of how many? Now say it as a percent.''
\end{itemize}
While circulating, pick the papers for the debrief --- one problem~13 that argues from 36 vs.\
20 (almost right: correct counts, wrong comparison), one that converts to 24\% vs.\ 40\%.
\textbf{Your formative read comes twice}: here, and again in the last eight minutes as students
start the homework.
\end{multicols}
\end{skillbox}

\boxguard
\begin{skillbox}[Debrief --- whole class (8 min)]{sky}
Pens down, papers up. \textbf{This debrief is spoken} --- there is no box at the end of the
notes to fill in, so it lives on the board and in the room.
\begin{multicols}{2}
\begin{enumerate}[leftmargin=1.3em, nosep, itemsep=3pt]
  \item \textbf{Open on the \emph{almost}-right problem~8 / 13 answer} you collected --- the
        one that says ``the older drivers text more, 36 to 20.'' Its counts are right; ask the
        room what question it answered. Then put the percents beside it.
  \item Put the four-line table from the debrief slide on the board: counts in one column,
        conditional percents in the other, for both contexts. The bigger count came from the
        bigger group both times.
  \item Take problem~6 back: \emph{``of 8+ hour sleepers''} and \emph{``of phone-in-room
        students''} use the same 48 and different denominators --- different questions.
  \item Close on the hedge: associated, yes; caused, not shown --- and with 40 students in one
        group, chance is not ruled out.
\end{enumerate}
\columnbreak
\textbf{Demand one sentence aloud} in AP form from a volunteer: \emph{``In this survey, 50\% of
students without a phone in the bedroom sleep 8+ hours, compared with 30\% of students with one,
so phone location and sleep are associated.''} Every part --- percents, both groups, context,
``associated'' not ``causes'' --- is what the free response scores.

\textbf{If the debrief runs short}, cut item~3. Item~1 is the only one that cannot be cut.
\textbf{Do not let the debrief eat the last eight minutes} --- the homework start is not optional
padding, it is your second formative read.
\end{multicols}
\end{skillbox}

\boxguard
\begin{skillbox}[AP Practice --- assigned, not scored]{goldbox}
One page on a 400-household grocery survey (location $\times$ grocery method, a $2\times3$
table): \textbf{MC~1} is the denominator trap --- ``of the households that use delivery'' is
$60/75 = 0.80$, and distractor (B) $0.240$ is the reverse conditional $60/250$; \textbf{MC~2} is
the crux, with distractor (A) arguing from raw counts (150 vs.\ 105) in groups of 250 and 150 ---
in rates the rural households actually shop in store \emph{more} (70\% vs.\ 60\%). The free
response computes the rural conditional distribution, judges association with the urban one,
and \textbf{part (c) is the spiral} to Lessons 1.2--1.3: a histogram needs a quantitative
variable; these are categories.

\textbf{Not scored --- the cover's score column reads \textbf{NA} for this page.} Go over it at
the start of Lesson~2.2 and hold the AP standard out loud: \emph{in context or it does not
count}. A part~(b) that says ``yes, the numbers are different'' without naming \emph{which}
percents for \emph{which} groups does not earn the point.
\end{skillbox}

\boxguard
\begin{skillbox}[Homework --- scored, due the first class after two study halls]{goldbox}
Five items in a third context --- a gym's 250 newest members, monthly vs.\ annual plan
$\times$ still attending after six months. \textbf{Part A:} A1 a joint relative frequency named
as joint; \textbf{A2 the denominator trap} (``of those who \emph{stopped}'' $= 90/120$).
\textbf{Part B:} B1 both conditional rates (40\% vs.\ 70\%); \textbf{B2 the crux restated} ---
the manager argues from 70 vs.\ 60 --- plus the causation hedge; B3 an AP-style multiple-choice
item (which result shows \emph{no} association) with a one-sentence justification.

\textbf{Start it in the last eight minutes of the period, in class and alone.}
\textbf{Scoring:} 10 points --- 2 per item. \textbf{B2 is the one to read closely}: it predicts
Lesson~2.2, where students must again describe a relationship in context and stop short of
cause.

\textbf{Paper homework is the default for this course} --- assign this page unless you have
decided otherwise. \textbf{DeltaMath override (occasional):} on the days you assign DeltaMath
instead, it replaces this page, you strike through the score cell on the cover, and this page
stays in the packet as extra practice. Never assign both.
\end{skillbox}

\boxguard
\begin{skillbox}[Watch For (while circulating)]{redbox}
\begin{itemize}[leftmargin=1.2em, nosep]
  \item \textbf{Comparing counts from groups of different sizes} (notes~8, 13; AP~MC2(A);
        HW~B2 --- the target misconception). Probe: \emph{``Out of how many? Say it as a
        percent.''}
  \item \textbf{The wrong denominator for a conditional} --- the grand total, or the other
        variable's total (notes~6; AP~MC1(B); HW~A2). Probe: \emph{``Read me the words after
        `of'. That group's total goes on the bottom.''}
  \item \textbf{Association read as causation} (notes~10, 14; HW~B2). Probe: \emph{``Did anyone
        assign the phones? Then what else could explain it?''}
  \item \textbf{``No association because the same category wins in both groups''} (AP~MC2(C)).
        Probe: \emph{``Do the bars match, or just their tallest piece?''}
  \item \textbf{Joint and marginal swapped} (notes~3, 11; HW~A1). Probe: \emph{``One cell or a
        whole row?''}
\end{itemize}
Cold-call prompts: \emph{``Of whom?''} \quad \emph{``Give me a pair of groups where the bigger
count is the smaller rate.''}
\end{skillbox}

\boxguard
\begin{skillbox}[Close \& Assign (8 min)]{goldbox}
\textbf{Students start the homework here, in class, alone and silent --- this is the individual
practice, and the first minutes of it are your best formative read.} Hand the launch first ---
five items on paper, scored, due the first class after two study halls --- and point out that
the AP Practice page before it is theirs to work and bring to review. Circulate:
\textbf{A2 and B2 tell you whether rows 2 and 3 landed.} Sort what you see into three piles ---
divided by 120 and argued B2 from 70\% vs.\ 40\% (ready); right numbers but no percents or no
hedge in B2 (ready, needs the language); divided by 150 or 250, or argued from 70 vs.\ 60 (the
misconception survived) --- and open Lesson~2.2 by reteaching row~3 with the warm-up's
\emph{more} vs.\ \emph{better} if that third pile ran past a third of the room. Then close on
what changed: \emph{``You walked in comparing counts. You walk out asking `of whom?' and comparing
percents --- and knowing that an association in a survey is a pattern, not a cause.''}
\textbf{Preview:} Lesson~2.2, \emph{Scatterplots and Correlation} --- the same question, ``are
these variables related?'', when both variables are numbers.
\end{skillbox}

% ── Teacher notes — one per component, in packet order (four: there is no activity) ──
\begin{teachernote}[Warm-Up]
Item~2 is the seed and the only item worth time. Give four minutes total. Expect half the room
to answer ``Section B'' to \emph{both} questions --- that is the misconception arriving early,
which is exactly what you want. Do not correct it with a lecture: have a student compute the two
percents aloud, then write \emph{more pass: B} and \emph{did better: A} side by side and leave
them on the board all period. Item~3's classification spirals Lesson~1.1; accept any sensible
``are these related?'' question and keep the phone/sleep one if a student offers it.
\end{teachernote}

\begin{teachernote}[Guided Notes \& Practice]
\textbf{This one document is the lesson --- pace it 24 / 10, then 8 for the debrief, then 8 for
the homework start.} Four rows, 14 problems. You work problems 1, 4, and 7; the class works the
rest of each grid; Guided Practice (11--14) is theirs with you asking. Row~1 is machinery and
must move --- six minutes. \textbf{Problem~6 is the trap} (the column-conditional $48/68$) and
\textbf{problem~8 is the crux}, in the last instruction row: the raw-count claim that runs
backwards once it is converted to rates. Spend the time there. If students need a reason
\emph{why} counts mislead, ask how many students are in each bar; the graph's labels (160 and 40)
are there for that.

\textbf{Guided Practice is where the pen changes hands} --- take answers, write what students
say, and push back rather than working it on the board while they watch. If you only get
through 12 and 13, that is the right trade.

\textbf{There is no independent practice set, and that is deliberate --- the homework is the
individual practice.} During the supervised start, read A2 and B2 over shoulders and sort into
the three piles in \emph{Close \& Assign}. If the third pile is more than a third of the room,
\textbf{open Lesson~2.2 by reteaching row~3}, and say so plainly.

\textbf{Debrief (8 min), spoken.} The almost-right crux answer first; it is the one move that
cannot be cut. \textbf{Do not let the debrief eat the last eight minutes.}
\end{teachernote}

\begin{teachernote}[AP Practice]
\textbf{This page is not required and not scored} --- the graded practice is the homework behind
it at the back of the packet. Go over it at the start of Lesson~2.2. \textbf{MC~1 is the one
worth reading closely}: a student who picks (B) 0.240 divided by the urban total when the question
restricted to delivery users --- the denominator trap in exam form. On the free response, part~(b)
earns credit only with both groups' percents for the same category, in context; part~(c) is the
spiral and should name the variable type, not just ``bar graphs are better.''
\end{teachernote}

\begin{teachernote}[Homework]
\textbf{Scored, 10 points (2 per item), due the first class after two study halls.} The eight
minutes at the end of the period are a \emph{start}: put students on A2 and B2 first while you
can still catch a wrong start. Scoring: A1 needs both the value and the word \emph{joint};
A2 earns nothing for $90/150$ or $90/250$; B2 needs both percents, a statement of association,
\emph{and} the observational-study hedge (1 point without the hedge); B3 needs (B) plus the
reason (equal conditional rates), not just the letter. \textbf{Paper is the default.} On the
occasions you assign DeltaMath in its place, strike the score cell on the cover and tell
students this page is now optional practice. Do not assign both.
\end{teachernote}
\end{document}
